\chapter{Giới hạn của Thước kẻ -- Compa và Tiên đề 6}

\section{Giới hạn đại số của hệ dựng hình Euclid}

Trong hình học Euclid phẳng, mọi phép dựng hình bằng thước kẻ không chia độ và compa đều quy về việc thực hiện hữu hạn các thao tác cơ bản sau:
\begin{enumerate}
    \item Tìm giao điểm của hai đường thẳng (tương đương giải hệ hai phương trình bậc nhất).
    \item Tìm giao điểm của một đường thẳng và một đường tròn (tương đương giải hệ gồm một phương trình bậc nhất và một phương trình bậc hai).
    \item Tìm giao điểm của hai đường tròn (tương đương giải hệ hai phương trình bậc hai).
\end{enumerate}

Về mặt đại số, nếu xuất phát từ một trường số cơ sở $\mathbb{Q}$, mỗi phép dựng bằng thước kẻ và compa chỉ cho phép ta mở rộng trường số bằng cách khai căn bậc hai. Nghĩa là, nếu một số $\alpha$ dựng được bằng thước kẻ và compa, thì bậc của trường mở rộng $[\mathbb{Q}(\alpha) : \mathbb{Q}]$ phải có dạng:
\begin{equation}
    [\mathbb{Q}(\alpha) : \mathbb{Q}] = 2^k \quad (k \in \mathbb{N}).
\end{equation}

Đối với bài toán nhân đôi khối lập phương, ta cần dựng đại lượng $x = \sqrt[3]{2}$, là nghiệm của đa thức bất khả quy $x^3 - 2 = 0$ trên $\mathbb{Q}$. Bậc mở rộng tương ứng là:
\begin{equation}
    [\mathbb{Q}(\sqrt[3]{2}) : \mathbb{Q}] = 3 \neq 2^k.
\end{equation}

Đối với bài toán chia ba góc, xét góc $\theta = 60^\circ$, việc chia ba góc này quy về việc tính $x = \cos(20^\circ)$. Dùng công thức nhân ba $\cos(3\theta) = 4\cos^3\theta - 3\cos\theta$, ta có phương trình:
\begin{equation}
    8x^3 - 6x - 1 = 0.
\end{equation}
Đa thức này bất khả quy trên $\mathbb{Q}$, do đó $[\mathbb{Q}(\cos 20^\circ) : \mathbb{Q}] = 3 \neq 2^k$.

Vì $3$ không phải là lũy thừa của $2$, cả hai bài toán đều hoàn toàn bất khả thi trong hệ dựng hình Euclid.

\section{Hệ 7 tiên đề Huzita--Hatori trong Origami}

Hệ thống toán học của Origami được hoàn thiện nhờ Humiaki Huzita (1989) và Koshiro Hatori (2001), bao gồm 7 tiên đề xác định cách tạo ra một nếp gấp thẳng duy nhất từ các điểm và đường thẳng cho trước:

\begin{enumerate}
    \item \textbf{Tiên đề 1:} Cho hai điểm $P_1, P_2$, có duy nhất một nếp gấp đi qua cả hai điểm (tương đương kẻ đường thẳng qua 2 điểm).
    \item \textbf{Tiên đề 2:} Cho hai điểm $P_1, P_2$, có duy nhất một nếp gấp đưa $P_1$ trùng vào $P_2$ (dựng đường trung trực).
    \item \textbf{Tiên đề 3:} Cho hai đường thẳng $L_1, L_2$, có thể gấp đưa $L_1$ trùng vào $L_2$ (dựng đường phân giác).
    \item \textbf{Tiên đề 4:} Cho điểm $P_1$ và đường thẳng $L_1$, có duy nhất một nếp gấp đi qua $P_1$ và vuông góc với $L_1$ (dựng đường vuông góc).
    \item \textbf{Tiên đề 5:} Cho hai điểm $P_1, P_2$ và đường thẳng $L_1$, có thể gấp đưa $P_1$ nằm trên $L_1$ sao cho nếp gấp đi qua $P_2$ (tương đương dựng giao điểm của đường tròn và đường thẳng).
    \item \textbf{Tiên đề 6:} Cho hai điểm $P_1, P_2$ và hai đường thẳng $L_1, L_2$, có thể tạo một nếp gấp đưa đồng thời $P_1$ lên $L_1$ và $P_2$ lên $L_2$.
    \item \textbf{Tiên đề 7:} Cho điểm $P_1$ và hai đường thẳng $L_1, L_2$, có thể tạo một nếp gấp đưa $P_1$ lên $L_1$ sao cho nếp gấp vuông góc với $L_2$.
\end{enumerate}

Các Tiên đề từ 1 đến 5 và Tiên đề 7 chỉ tương đương với các phép toán giải phương trình bậc nhất và bậc hai. Riêng \textbf{Tiên đề 6} mang lại bước nhảy vọt về năng lực dựng hình.

\section{Bản chất hình học của Tiên đề 6: Tiếp tuyến chung của hai Parabol}

Xét thao tác gấp một điểm $P$ lên một đường thẳng $L$ bởi nếp gấp $T$:
\begin{itemize}
    \item Nếp gấp $T$ chính là đường trung trực của đoạn nối điểm $P$ với ảnh $P'$ của nó trên $L$.
    \item Theo định nghĩa quỹ tích hình học, đường thẳng $T$ luôn tiếp xúc với một parabol có tiêu điểm (focus) là $P$ và đường chuẩn (directrix) là $L$.
\end{itemize}

Khi thực hiện Tiên đề 6:
\begin{itemize}
    \item Thao tác đưa $P_1$ lên $L_1$ đòi hỏi nếp gấp $T$ phải tiếp xúc với Parabol $\mathcal{P}_1$ (tiêu điểm $P_1$, đường chuẩn $L_1$).
    \item Thao tác đồng thời đưa $P_2$ lên $L_2$ đòi hỏi $T$ tiếp xúc với Parabol $\mathcal{P}_2$ (tiêu điểm $P_2$, đường chuẩn $L_2$).
\end{itemize}

Do đó, nếp gấp $T$ chính là \textbf{tiếp tuyến chung} của hai đường parabol $\mathcal{P}_1$ và $\mathcal{P}_2$.

\begin{figure}[H]
    \centering
    \begin{tikzpicture}[scale=1.1, >=stealth]
        % --- Khung nền tờ giấy Origami ---
        \fill[gray!5] (-2,-2) rectangle (5,5);
        \draw[thick, gray!60] (-2,-2) rectangle (5,5);

        % --- Đường chuẩn L1 và L2 ---
        \draw[very thick, blue!80!black] (-1, -1.8) -- (-1, 4.8) node[above] {\small $L_1$ (chuẩn của $\mathcal{P}_1$)};
        \draw[very thick, red!80!black] (-1.8, -1) -- (4.8, -1) node[right] {\small $L_2$ (chuẩn của $\mathcal{P}_2$)};

        % --- Tiêu điểm P1 và P2 ---
        \coordinate (P1) at (1, 2);
        \coordinate (P2) at (2, 1);
        \fill[blue!90!black] (P1) circle (2.2pt) node[above left] {$P_1$};
        \fill[red!90!black] (P2) circle (2.2pt) node[below right] {$P_2$};

        % --- Hai đường Parabol P1 và P2 ---
        % Parabol P1: Tiêu điểm (1,2), đường chuẩn x = -1 -> Đỉnh (0,2), pt: x = (y-2)^2 / 4
        \draw[thick, blue!65, domain=-0.8:4.8, smooth, variable=\y] 
            plot ({((\y-2)^2)/4}, \y);
        \node[blue!80!black, font=\footnotesize] at (3.5, 4.4) {$\mathcal{P}_1$ (tiêu điểm $P_1$, chuẩn $L_1$)};

        % Parabol P2: Tiêu điểm (2,1), đường chuẩn y = -1 -> Đỉnh (2,0), pt: y = (x-2)^2 / 4
        \draw[thick, red!65, domain=-0.8:4.8, smooth, variable=\x] 
            plot (\x, {((\x-2)^2)/4});
        \node[red!80!black, font=\footnotesize] at (4.4, 2.0) {$\mathcal{P}_2$ (tiêu điểm $P_2$, chuẩn $L_2$)};

        % --- Nếp gấp T: Tiếp tuyến chung đối xứng x + y = 1 (y = -x + 1) ---
        % Tiếp xúc P1 tại (1, 0) và P2 tại (0, 1)
        \draw[very thick, teal, dashed, domain=-1.6:2.6] 
            plot (\x, {-\x + 1});
        \node[teal!90!black, font=\bfseries, rotate=-45] at (2.2, -0.8) {Nếp gấp $T$ (tiếp tuyến chung)};

        % --- Các điểm tiếp xúc của T với 2 Parabol ---
        \coordinate (Touch1) at (1, 0);
        \coordinate (Touch2) at (0, 1);
        \fill[teal] (Touch1) circle (2pt) node[below left, font=\scriptsize] {Tiếp xúc $\mathcal{P}_1$};
        \fill[teal] (Touch2) circle (2pt) node[above right, font=\scriptsize] {Tiếp xúc $\mathcal{P}_2$};

        % --- Minh họa điểm ảnh P1' in L1 và P2' in L2 qua nếp gấp ---
        \coordinate (P1_prime) at (-1, 0);
        \coordinate (P2_prime) at (0, -1);
        \fill[blue!50!black] (P1_prime) circle (1.8pt) node[left] {$P'_1 \in L_1$};
        \fill[red!50!black] (P2_prime) circle (1.8pt) node[below] {$P'_2 \in L_2$};

        % Đường gióng phản xạ vuông góc với nếp gấp T
        \draw[densely dotted, blue!70] (P1) -- (P1_prime);
        \draw[densely dotted, red!70] (P2) -- (P2_prime);
    \end{tikzpicture}
    \caption{Cơ sở hình học của Tiên đề 6: Nếp gấp $T$ đưa đồng thời $P_1 \to L_1$ và $P_2 \to L_2$ tiếp xúc với $\mathcal{P}_1$ tại $(1,0)$ và $\mathcal{P}_2$ tại $(0,1)$, trở thành tiếp tuyến chung của hai parabol.}
    \label{fig:axiom6_tangent_steps}
\end{figure}

\section{Ý nghĩa đại số: Giải phương trình bậc ba}

Giả sử nếp gấp $T$ có hệ số góc là $m$. Điều kiện để một đường thẳng tiếp xúc với một parabol bậc hai dẫn đến một phương trình liên hệ giữa $m$ và các tham số của parabol.

Khi xét đồng thời hai parabol có trục không song song, việc tìm tiếp tuyến chung dẫn đến việc giải một hệ phương trình mà sau khi khử tham số sẽ thu được một \textbf{phương trình bậc ba} đối với hệ số góc $m$:
\begin{equation}
    a_3 m^3 + a_2 m^2 + a_1 m + a_0 = 0 \quad (a_3 \neq 0).
\end{equation}

Hai parabol có thể có tối đa 3 tiếp tuyến chung thực, tương ứng với tối đa 3 nghiệm thực của phương trình bậc ba. Thao tác gập giấy vật lý trong Tiên đề 6 cho phép ta trượt và xoay tờ giấy liên tục cho đến khi hai điểm chạm đồng thời vào hai đường thẳng, tức là trực tiếp tìm ra nghiệm thực của phương trình bậc ba mà không cần công thức Cardano phức tạp.


Vì trường số dựng được bởi Origami có thể mở rộng bậc 3, ta hoàn toàn có thể xây dựng các trường mở rộng có dạng:
\begin{equation}
    [\mathbb{K} : \mathbb{Q}] = 2^a \cdot 3^b \quad (a, b \in \mathbb{N}).
\end{equation}
Đây chính là lý do toán học giải thích vì sao Origami giải quyết được triệt để cả bài toán chia ba góc và nhân đôi khối lập phương.